\documentclass[12pt,a4paper,onecolumn]{article}

% Encoding and language
\usepackage[utf8]{inputenc}
\usepackage[T1]{fontenc}
\usepackage[english]{babel}

% Font
\usepackage{times}

% Set margins
\usepackage[a4paper, top=2cm, bottom=2cm, left=2.5cm, right=2.5cm]{geometry}

% Tabulation
\usepackage{tabto}

% Line spacing
\usepackage{setspace}
\onehalfspacing

% Page numbers
\usepackage{fancyhdr}
\pagestyle{fancy}
\fancyhf{} % Clear default header/footer
\rfoot{\thepage} % Right-aligned page numbers in footer

% Remove page number from cover/abstract/etc
\usepackage{titlesec}
\newcommand{\noPageNumber}{\thispagestyle{empty}}

% Prevent indentation
\setlength{\parindent}{0pt}
\setlength{\parskip}{0.5em}

% Figure management
\usepackage{graphicx}
\usepackage{caption}
\usepackage{subcaption}

% Table
\usepackage{booktabs, tabularx}
\renewcommand\tabularxcolumn[1]{m{#1}}

\usepackage{pgfplots}
\pgfplotsset{compat=1.16}
\def\mathdefault#1{#1}
\everymath=\expandafter{\the\everymath\displaystyle}

% Subfigure management
\usepackage{layouts}
%\usepackage{floatrow}
%\usepackage[label font=bf,labelformat=simple]{subfig}% <-- changed
%\floatsetup[figure]{style=plain,subcapbesideposition=top}


% Optional: headings formatting
\titleformat{\section}{\large\bfseries}{\thesection}{1em}{}
\titleformat{\subsection}{\normalsize\bfseries}{\thesubsection}{1em}{}

% Optional: remove widow/orphan lines
\widowpenalty=10000
\clubpenalty=10000

% Mathematics
\usepackage{amsmath,mathtools,amssymb} 

\begin{document}
 %% textwidth in cm: \printinunitsof{in}\prntlen{\textwidth}

\begin{table}[htb]
\centering
\footnotesize
\begin{tabularx}{\linewidth}{c >{\raggedright\arraybackslash}X c c c}
\toprule
Parameter & Definition & Unit & Par A & Par B \\

\midrule
\textit{n} & Number of host LCT subcompartments. Equals to $1/CV_H^2$ with $CV_H$ the
coefficient of variation of host maturation time & unit-less & \textit{varies} & \textit{varies} \\

\textit{m} & Number of parasitoid LCT subcompartments. Equals to $1/CV_P^2$ with $CV_P$ the
coefficient of variation of parasitoid maturation time & unit-less & \textit{varies} & \textit{varies} \\

$m_{H1}$ & Maturation rate of H1. Average maturation time  $T_{H1}$, can be defined as $1/m_{H1}$ & $t^{-1}$ & 1  & 1 \\

$\delta_{H1}$ & Background mortality of H1 & $t^{-1}$ & 0 & 0.00425 \\
$\delta_{H2}$ & Background mortality of H2. Host adult longevity, $T_{H2}$ can be defined as $1/\delta_{H2}$ & $t^{-1}$ &  0.2 & 0.2  \\
$^\rho$ & Average number of offspring of one adult during its lifetime. Equals to $r/\delta_{H2}$ with r the adult fecundity rate. & unit-less & 33 & 105 \\

$m_{P1}$ & Maturation rate of P1. Average maturation time  $T_{P1}$, can be defined as $1/m_{P1}$ & $t^{-1}$ & 1 & 0.05 \\
$\delta_{P1}$ & Background mortality of P1 & $t^{-1}$ & 1 & 0.1  \\
$\delta_{P2}$ & Background mortality of P2 & $t^{-1}$ & 8 & 0.5  \\
$a$ & Rate of host discovery by a single parasitoid & $t^{-1}$ & 1 & 0.03  \\

$\theta_H$ & Rate parameter of Erlang maturation time distribution. Set equal to $n \cdot m_{H1}$ (Appendix A.1.3) & $t^{-1}$ & \textit{varies}  & \textit{varies}   \\




\bottomrule
\end{tabularx}
\caption{Parameters and values of Model I (Murdoch 1987). Par A : parameter set inspired from those in Murdoch (1987). Par B : parameter set inspired from measurements in Sait et al. (1994) and Harvey et al. (1994) on the \textit{Plodia-Venturia} host-parasitoid system.}
\end{table}




























\begin{table}[htb]
\centering
\footnotesize
\begin{tabularx}{\linewidth}{c >{\raggedright\arraybackslash}X c c c}
\toprule
Parameter & Definition & Unit & Value \\

\midrule
$n_1$ & Number of H1 LCT subcompartments. Equals to $1/CV_{H1}^2$ with $CV_{H1}$ the
coefficient of variation of H1 maturation time & unit-less & \textit{varies}  \\

$n_2$ & Number of H2 LCT subcompartments. Equals to $1/CV_{H2}^2$ with $CV_{H2}$ the
coefficient of variation of H2 maturation time & unit-less & \textit{varies}  \\

$n_3$ & Number of H3 LCT subcompartments. Equals to $1/CV_{H3}^2$ with $CV_{H3}$ the
coefficient of variation of H3 maturation time & unit-less & \textit{varies}  \\

\textit{m} & Number of parasite LCT subcompartments. Equals to $1/CV_P^2$ with $CV_P$ the
coefficient of variation of parasitoid maturation time & unit-less & \textit{varies}  \\

$m_{H1}$ & Maturation rate of H1 & $t^{-1}$ & 1/3 \\
$m_{H2}$ & Maturation rate of H2 & $t^{-1}$ & 1/2  \\
$m_{H3}$ & Maturation rate of H3 & $t^{-1}$ & 1/3  \\

$\delta_{H1}$ & Background mortality of H1 & $t^{-1}$ & 0  \\
$\delta_{H2}$ & Background mortality of H2 & $t^{-1}$ & 0  \\
$\delta_{H3}$ & Background mortality of H3 & $t^{-1}$ & 0  \\
$\delta_{H4}$ & Background mortality of H4. Host adult longevity, $T_{H4}$ can be defined as $1/\delta_{H4}$ & $t^{-1}$ &  0.2 \\
$^\rho$ & Average number of offspring of one adult during its lifetime. Equals to $r/\delta_{H4}$ with r the adult fecundity rate. & unit-less & 15  \\

$m_{P1}$ & Maturation rate of P1 & $t^{-1}$ & \textit{varies}  \\
$\delta_{P1}$ & Background mortality of P1 & $t^{-1}$ & 0   \\
$\delta_{P2}$ & Background mortality of P2 & $t^{-1}$ & 2.5   \\
$a$ & Rate of host discovery by a single parasitoid & $t^{-1}$ & 0.01   \\
$k$ & Parasitoid aggregation or clumping. Under specific conditions, one has the equivalence $k \approx n/\hat \tau$ with $\hat \tau$ the average maturation time (Appendix A.3). & unit-less & \textit{varies}   \\

$\theta_X$ & Rate parameter of Erlang maturation time distribution of X : H1, H2, H3. Set equal to $n \cdot m_{X}$ & $t^{-1}$ & \textit{varies}     \\

$\textit{ratioT}$ & Ratio between parasitoid and host generation time $= \frac{ \frac{1}{m_{H1}} + \frac{1}{m_{H2}}+ \frac{1}{m_{H3}} + \frac{1}{\delta_{H4}} }{\frac{1}{m_{P1}} + \frac{1}{\delta_{P2}} }$ & unit-less & \textit{varies}  \\

\bottomrule
\end{tabularx}
\caption{Parameters and values of Model II (Godfray and Hassel 1989)}
\end{table}


\end{document}